% Laje, caso "paredes" (problems/screened_poisson/slab_cases.py), Metodo 1 sem treino:
% g0 = tau* em S contra o g* livre que minimiza o erro de LAT.
% tools/camps_gstar/fit_g_slab.py (MESH 49,25,5, EPS 2; .119, 2026-10-16);
% figuras: scripts/slab_fig.py (data/slab/gstar_slab.npz). So LAT: sem ECG na laje.

\section{Laje: $g$ ótimo $\times$ verdade}

\begin{frame}{Laje, caso paredes: $g_0 = \tau^*$ em $S$ $\times$ $g^*$ ótimo}
  \footnotesize
  \begin{itemize}\setlength{\itemsep}{0.1em}
    \item Laje $2 \times 1 \times 0{,}15$, malha $49 \times 25 \times 5$ (6125 nós), fibra em $x$, $c = (0{,}052;\ 0{,}017)$ cm/ms.
    \item $S$ = paredes $y = 0$ e $x = 0$ (365 nós). Fonte verdadeira: toda a parede $y = 0$ em $t = 0$ (onda plana); verdade exata.
    \item Método 1 sem treino: $\varepsilon = 2$ ms, massa agrupada. $g_0 = \tau^*$ em $S$;
      $g^*$ = $g$ livre em $S$ que minimiza o erro de LAT (Adam, 3000 passos, a partir de $g_0$).
  \end{itemize}
  \medskip
  \centering
  \scriptsize
  \setlength{\tabcolsep}{5pt}
  \begin{tabular}{@{}lcccccccc@{}}
    \toprule
    & \shortstack{perda\\de LAT} & \shortstack{LAT MAE\\(ms)} & \shortstack{LAT rms\\(ms)} & corr & eik & bacias & \shortstack{dispersão\\na fonte (ms)} & \shortstack{nós\\$u \leq 0$}\\
    \midrule
    verdade & -- & -- & -- & -- & 0,000 & 1 & 0,00 & --\\
    $g_0 = \tau^*$ em $S$ & $2{,}6 \cdot 10^{-4}$ & 0,214 & 0,286 & 0,99987 & 0,005 & 2 & 1,26 & 0\%\\
    $g^*$ (ajuste ao LAT) & $8{,}2 \cdot 10^{-7}$ & 0,017 & 0,028 & 1,00000 & 0,004 & 15 & 0,19 & 0\%\\
    \bottomrule
  \end{tabular}\par
  \medskip
  {\tiny Só LAT (a laje não tem ECG aqui). MAE e rms com as médias removidas. eik: RMS de $\|\U\nabla\varphi\| - 1$.
  bacias: mínimos locais do mapa. dispersão na fonte: $\max - \min$ de $\varphi$ na parede $y = 0$ (0 na verdade).}
\end{frame}

\begin{frame}{Laje: mapas e erro de LAT}
  \centering
  \includegraphics[width=\linewidth]{laje_gstar_mapas.png}\par
  {\scriptsize Plano de meia altura. Amarelo: parte de $S$ que é fonte; roxo: parte de $S$ que não é fonte.}
\end{frame}

\begin{frame}{Laje: o que o ajuste muda em $g$}
  \centering
  \includegraphics[width=0.92\linewidth]{laje_gstar_g.png}\par
  \medskip
  \small
  \begin{itemize}\setlength{\itemsep}{0.2em}
    \item Na parede $y = 0$ (a fonte), $g^*$ fica a 0,5--1 ms do nó mais cedo; até 3,6 ms no canto com a parede $x = 0$.
    \item Na parede $x = 0$ (em $S$, mas não é fonte), $g^*$ fica 5--8 ms depois de $\tau^*$: com $g_0 = \tau^*$ esses nós injetam fonte quando a frente passa e adiantam o mapa junto à parede (até 1,1 ms).
  \end{itemize}
\end{frame}
